\documentclass[11pt]{article}

\usepackage[a4paper,margin=1in]{geometry}
\usepackage{amsmath,amssymb,amsthm}
\usepackage{enumitem}
\usepackage{microtype}
\usepackage{parskip}

% -------------------------------------------------------------------
% Environments
% -------------------------------------------------------------------
\theoremstyle{definition}
\newtheorem*{definition}{Definition}
\newtheorem*{question}{Question}
\newtheorem*{example}{Example}

% -------------------------------------------------------------------
% Notation
% -------------------------------------------------------------------
% The congruence class a (mod m): \cls{a}{m}
\newcommand{\cls}[2]{#1 \pmod{#2}}

% A problem cell: \cell{<number>}{<title>}{<points>}
\newcommand{\cell}[3]{%
  \section*{Cell #1: #2\hfill{\normalsize\mdseries #3}}}

\title{Disjoint Congruence Classes\\[0.4em]\large Problem Description}
\author{}
\date{}

\begin{document}
\maketitle

% ===================================================================
\section*{Overview}
% ===================================================================

If congruence classes are pairwise disjoint, must two of their moduli
share a large common factor?

\begin{definition}[Congruence class]
For integers $a$ and $m \geq 1$, write
\[
  \cls{a}{m} = \{\, a + mt : t \in \mathbb{Z} \,\}
\]
for the \emph{congruence class} of $a$ modulo $m$.
\end{definition}

\begin{definition}[Pairwise disjoint]
A finite family of classes
\[
  \cls{a_1}{m_1}, \; \ldots, \; \cls{a_k}{m_k}
\]
is \emph{pairwise disjoint} if
\[
  \bigl(\cls{a_i}{m_i}\bigr) \cap \bigl(\cls{a_j}{m_j}\bigr)
  = \varnothing
  \qquad \text{whenever } i < j.
\]
Nothing is assumed about the moduli beyond $m_i \geq 1$: they may
repeat, and the residues $a_i$ may repeat as well.
\end{definition}

By the Chinese remainder theorem, two classes meet if and only if the
congruences
\[
  x \equiv a_i \pmod{m_i},
  \qquad
  x \equiv a_j \pmod{m_j}
\]
are compatible, that is,
\begin{equation}
  \bigl(\cls{a_i}{m_i}\bigr) \cap \bigl(\cls{a_j}{m_j}\bigr)
  \neq \varnothing
  \quad \Longleftrightarrow \quad
  \gcd(m_i, m_j) \mid (a_i - a_j).
  \tag{$*$}\label{eq:intersection}
\end{equation}

The question of this column is how large the moduli of a disjoint
family must be forced to overlap in the gcd sense.

\begin{question}
Let $k \geq 2$ and let
\[
  \cls{a_1}{m_1}, \; \ldots, \; \cls{a_k}{m_k}
\]
be pairwise disjoint. Must there be a pair $i < j$ with
\[
  \gcd(m_i, m_j) \geq k?
\]
\end{question}

This is believed to be true for every $k$. It is sharp: the $k$
classes
\[
  \cls{1}{k}, \; \cls{2}{k}, \; \ldots, \; \cls{k}{k}
\]
are pairwise disjoint---their differences have absolute value less
than $k$, so $k$ divides none of them---and every pairwise gcd equals
exactly $k$. So $k$ could not be replaced by $k + 1$ in the statement.

\begin{example}
The three classes
\[
  \cls{0}{2}, \qquad \cls{1}{4}, \qquad \cls{3}{8}
\]
are pairwise disjoint by \eqref{eq:intersection}, and their pairwise
gcds are $2, 2, 4$; the largest is $4 \geq 3$, as the question
predicts for $k = 3$.

By contrast, $\cls{0}{2}$ and $\cls{0}{3}$ have coprime moduli and do
meet, at $0$.
\end{example}

\subsection*{Trivial Bounds}

Two observations are immediate and you should assume they are already
on the table; presenting either as progress counts for nothing.

\begin{enumerate}[leftmargin=*]
  \item If $\gcd(m_i, m_j) = 1$, then by \eqref{eq:intersection} the
    two classes meet. Hence in any pairwise disjoint family every
    pairwise gcd is at least $2$. This answers the question for
    $k = 2$, and for a family of size $k$ that fails the question it
    confines every pairwise gcd to the range
    \[
      2 \leq \gcd(m_i, m_j) \leq k - 1.
    \]
  \item The class $\cls{a}{m}$ has density $1/m$, so a pairwise
    disjoint family satisfies
    \[
      \sum_{i=1}^{k} \frac{1}{m_i} \leq 1.
    \]
    In particular, some modulus satisfies $m_i \geq k$.
\end{enumerate}

Observation~2 bounds a single modulus from below; it says nothing
about any gcd, and the two together settle no size $k \geq 3$. Every
cell below needs a genuinely new argument.

% ===================================================================
\section*{What to Hand In}
% ===================================================================

For each cell you attempt, hand in a written proof.

\begin{itemize}[leftmargin=*]
  \item \textbf{Computation.}
    You may use a computer to explore. A proof may rely on a
    computation only if you include the code, it runs in under
    10 minutes on a laptop, and the computation is exhaustive over a
    finite set that your own argument has reduced the problem to.
  \item \textbf{Exhaustiveness certificate.}
    A computational cell is accepted only with an
    \emph{exhaustiveness certificate}, meaning all five:
    % The source says ``all four'' but lists five requirements.
    \begin{enumerate}[leftmargin=*]
      \item An exact description of the finite set your search
        enumerates, together with the proof that nothing outside it
        can be a counterexample.
      \item Every pruning rule you use, stated precisely, each with
        the proof that it discards only objects that cannot be
        completed to a counterexample.
      \item The node count of the search at each size you claim, and
        the wall clock.
      \item A rerun that reproduces those counts.
      \item If anything at all survives your pruning at a size you
        claim, a separate decision for \textbf{every} survivor---not a
        sample---together with, for each one, the \textbf{smallest}
        part of it that already forces the decision, and your proof
        that no smaller part does. A survivor you cannot decide means
        the size is not certified.
    \end{enumerate}
    A run that reports ``I searched and found nothing'' without 1--5
    is not a certificate and scores nothing.
  \item \textbf{Unfinished searches.}
    If a search does not finish at some size, report it as unfinished
    and do not claim that size; a partial search is a legitimate
    partial result and should be handed in as one, with the counts it
    reached.
  \item \textbf{Boundaries.}
    Where a cell asks you to determine a boundary, state the exact
    value you claim. Claim only what you have verified: an unfinished
    size, or one certified under an assumption you have not verified,
    must be reported as such.
  \item \textbf{Status.}
    Say clearly which cells you consider solved and which are
    partial.
  \item \textbf{Citations.}
    Citing a published result for the statement you are asked to
    prove does not count as a solution; a proof written out in full
    does, whatever its source. If you rely on a published argument,
    you are responsible for it: check it yourself, and report it if
    it does not hold up.
\end{itemize}

% ===================================================================
\cell{1}{Three Classes}{}
% ===================================================================

\emph{The statement of this cell is not visible in the provided
screenshots.}

% ===================================================================
\cell{2}{Four Classes}{2 points}
% ===================================================================

The same question for four classes. Prove the statement for $k = 4$.

% ===================================================================
\cell{3}{Every $k$ up to 8}{3 points}
% ===================================================================

A first range of sizes: after C1 and C2, the sizes $k = 5, 6, 7$ and
$8$ remain. Prove the statement for every $k \leq 8$.

% ===================================================================
\cell{4}{Every $k$ up to 12}{5 points}
% ===================================================================

A longer range, and a precise account of your method at the sizes
just beyond it. Prove the statement for every $k \leq 12$.

Then, for each $k$ from $9$ to $16$ in turn, report exactly what your
method leaves undecided at that $k$: if nothing, say so and prove it;
if something, exhibit it in full.

If that disagrees with any source you consulted, say which of the two
is right and why. An answer that reports agreement with a source it
did not test will be marked wrong.

% ===================================================================
\cell{5}{The Certified Boundary}{8 points}
% ===================================================================

The main certification cell: an initial range of sizes as long as you
can make it, plus two isolated sizes further out.

Determine the largest $k$ for which you can certify the statement for
all sizes up to and including $k$, and certify it. The range you
claim must be contiguous, and if your argument at a given size
assumes that all smaller sizes have already been settled you must say
so.

Then decide the two isolated sizes $k = 24$ and $k = 30$: show that
neither is the least size at which the statement can fail.

For this cell the exhaustiveness certificate of the hand-in rules is
not enough on its own. Add:

\begin{enumerate}[label=(\roman*), leftmargin=*]
  \item \textbf{A demonstration that your method is not vacuous.}
    Construct a case in which an admissible configuration genuinely
    exists, run your machinery on it unmodified, and show it returns
    that configuration. A method that reports ``nothing survives'' at
    every size it is pointed at is indistinguishable from a method
    with a bug, and will be graded as one.
  \item \textbf{For every object your pruning leaves undecided, at
    every $k$ you claim}---not a sample---an individual decision, plus
    the smallest part of that object which already forces the
    decision, plus a proof that no smaller part does.
  \item \textbf{The exact list}, not merely the count, of what
    survives at each $k$.
  \item \textbf{Every pruning rule you use beyond those you have
    proved, proved.}
    If your search needs a rule you invented to finish a size, that
    rule is part of your claim: state it, prove it, and show the
    survivor list is unchanged when you switch it off. A size that
    only closes with an unproved rule is not certified.
  \item \textbf{A second implementation, written independently of
    your first, that differs in method}---not the same algorithm
    twice---and the two survivor lists compared elementwise at every
    $k$ you claim. Report any difference rather than reconciling it
    silently: a disagreement means one of them is wrong, and finding
    which is part of the cell.
\end{enumerate}

% ===================================================================
\cell{6}{Beyond the Boundary}{13 points}
% ===================================================================

\textbf{Open question.}

Three directions beyond the certified range; any one of them counts.
Any of the following.

\begin{enumerate}[label=(\alph*), leftmargin=*]
  \item Decide a size $k \geq 25$.
  \item \textbf{The asymptotic form.}
    It is known that a pairwise disjoint family of size $k$ always
    has a pair with
    \[
      \gcd(m_i, m_j)
      \geq
      k \cdot \exp\left( -(2 + o(1)) \sqrt{\frac{\log k}{\log\log k}} \right),
    \]
    which is $k^{1 - o(1)}$ but not linear in $k$. Prove the
    statement in full, or prove the weaker bound
    \[
      \gcd(m_i, m_j) \geq ck
    \]
    for some absolute constant $c > 0$, or improve the exponential
    factor above.
  \item \textbf{The group form.}
    Let $G$ be a group, let $G_1, \ldots, G_k$ be subgroups of finite
    index
    \[
      n_i = [G : G_i],
    \]
    and let $x_1 G_1, \ldots, x_k G_k$ be pairwise disjoint cosets.
    Is there a pair $i < j$ with
    \[
      \gcd(n_i, n_j) \geq k?
    \]
    This is known for $k \leq 5$ and open for every $k \geq 6$;
    settling $k = 6$ counts as progress.
\end{enumerate}

\end{document}
